% Section VII: Conclusion (target 0.15 pages). Numbers repeat Section VI; no new claims.
\section{Conclusion}
\label{sec:conclusion}
This paper presented a task-space imitation learning system for lunchbox
packing in which a seven-joint leader arm teaches a six-joint collaborative
robot. A clutch-relative mapping in cylindrical coordinates and one
command-shaping stage serve
both teleoperation and autonomous execution, so the recorded action is the
target the robot executes. Trained with ACT on 140 demonstrations, one policy
completed 22 of 30 packing trials over six layouts and placed 73 of 90 items
correctly, with no item released into a wrong compartment and a median
completion time of 66.6~s. The failures arose in grasping, mainly retries
that drifted outward to the workspace bound and items flicked or pushed away
by the fingertips. Extending the data beyond the limits of
Section~\ref{sec:limitations}, and residual reinforcement learning, remain
future work.
